\section{前提を確かめる：バイアス}

  \subsection{集め方の偏りと人数の限界}

    前章では, 原作の有無による差が本物の効果なのかを調べるために, 条件をそろえて比べる方法を考えた.
    くじで分けられないデータから差の理由を言うには, 「ほかの条件に大きな違いはないはずだ」という仮定に頼るしかなかった.

    だが, その話の足元には, もっと手前で置いたきりの前提が残っている.
    標本を取り直せば平均が散らばるという話をしたときから, 手元の80本は「対象作品全体から大きな偏りなく抽出された標本である」として扱ってきた.
    区間を作っても, 差を検定しても, 条件をそろえても, 出てきた結果を作品全体の話として言えるのは, この前提があってこそだった.
    もし集めてきたデータが最初から偏っていたとしたら, ここまでの計算はどうなってしまうだろうか.

    極端な場面を考えてみよう.
    たとえば日本に住む有権者全体の意見を知りたいのに, 平日の昼に渋谷駅にいた人だけに話を聞いたとする.
    平日の昼間に渋谷駅を歩いている人たちと, 有権者全体とでは, 年齢層も職業も日頃の暮らしもまるで違うはずだ.
    この方法で10万人の意見を集めたとしよう.
    10万人といえば途方もない規模だが, いくら人数を積み上げても, 平日の昼に渋谷駅にいなかった人たちの声は最初から最後までゼロのままだ.

    人数を増やすと, 標準誤差の分母にある $\sqrt{n}$ が大きくなるため, 計算される信頼区間はどこまでも狭くなっていく.
    一見すると, 推定の精度がどんどん上がっているように思えるかもしれない.
    しかし, 区間の中心である標本平均が近づいていく先は, 知りたかった有権者全体の平均ではなく, 「平日の昼に渋谷駅にいた人たち」の平均でしかない (\cref{fig:ch07-biased-n}).
    中心そのものがずれているのだから, 区間をどれほど狭く絞り込んでも, 本当に知りたい値から外れた場所を精密に指し示し続けるだけになってしまう.

    \begin{figure}[htbp]
      \centering
      % 第7章 図7-1: 偏った集め方のまま人数を増やしたときの信頼区間 (模式図, 目盛りなし)
% 区間の半分の長さは 1/sqrt(n) に比例させた: 100人 3.0, 1000人 0.95, 10万人 0.095
\begin{tikzpicture}[x=0.75cm, y=1cm, font=\footnotesize]
  % 知りたい値と, 偏った集まりの平均
  \draw[densely dashed] (0,-0.3) -- (0,3.3);
  \draw[densely dotted, thick] (6,-0.3) -- (6,3.3);
  \node[above, align=center] at (0,3.3) {有権者全体の平均};
  \node[above, align=center] at (6,3.3) {平日の昼に渋谷駅に\\いた人たちの平均};
  % 区間
  \foreach \y/\c/\h/\lab in {2.4/5.5/3.0/100人, 1.4/6.3/0.95/1000人, 0.4/5.97/0.095/10万人} {
    \draw[line width=2.2pt, black!55] ({\c-\h},\y) -- ({\c+\h},\y);
    \draw[thick] ({\c-\h},\y-0.14) -- ({\c-\h},\y+0.14) ({\c+\h},\y-0.14) -- ({\c+\h},\y+0.14);
    \fill (\c,\y) circle (2.2pt);
    \node[left] at (-0.6,\y) {\lab};
  }
\end{tikzpicture}

      \caption{平日の昼に渋谷駅にいた人だけに話を聞く方法で, 聞く人数を100人, 1000人, 10万人と増やしたときの信頼区間の模式図 (目盛りは付けていない). 黒い点は区間の中心の標本平均. 破線は有権者全体の平均, 点線は平日の昼に渋谷駅にいた人たちの平均の位置を示す. 区間の長さは, 人数の平方根に反比例するように描いた.}
      \label{fig:ch07-biased-n}
    \end{figure}

    同じ手順で作った区間には100回に95回ほど母平均が入る, という約束が成り立つのは, 標本が知りたい母集団全体から偏りなく選ばれているときだけだ.
    もし集め方が最初から偏っていれば, 同じ調査をどれほど繰り返したところで, 100回に95回入るのは母集団全体の平均ではなく, 偏って集められた集まりの平均である.
    このように, 標本を集める手順そのものに偏りがあることを\textbf{サンプリングバイアス}と呼ぶ.

    選挙の前に, ネットで「誰に投票しますか」と尋ねる投票が行われ, 10万票が集まって, 候補Aが7割の票を取ったとする.
    同じころ, 有権者から無作為に選んだ1000人に聞いた調査では, 候補Aは32\%と出ていた.
    実際の選挙では, 候補Aの得票は30\%だった.
    当たっていたのは, 100分の1の人数しか聞いていない調査のほうだ.

    1000人に聞けば, 世論調査の支持率で計算したとおり, $\pm 3$ポイントほどの幅がつく.
    32\%の上下3ポイントには, 30\%がちゃんと入っている.
    同じ計算を10万票に当てはめると, 幅は $\pm 0.3$ポイントほどまで縮む.
    その狭い区間が, 実際の30\%から40ポイントも離れたところにあった.

    ネットで票を入れたのは, 自分から投票しに来た人たちだけだった.
    票数が増えるほど, その人たちの意見を精密に指すようになるだけで, 有権者全体には近づかない.
    \cref{fig:ch07-biased-n} でいえば, 10万票の7割は, 点線の位置を短い区間で指していたことになる.

    1000人の調査とネットの投票を分けたのは, 人数ではなく, 誰の声が集まったかだった.
    偏りは, 聞く相手を選ぶ前から入り込んでいることもある.
    映画80本にも, それが起こりうる.
    80本を選んだ元の一覧に並んでいるのは, いまそのサービスで配信されている作品だ.
    公開から30日間ほとんど見られず, その後に配信を終えた作品があったとしても, 一覧にはもう載っていない.
    母集団をこれまでどおり配信されている作品全体とするなら, これは偏りではない.
    消えた作品は, はじめから母集団に入っていないからだ.
    だが, 知りたいのが「映画は配信でどれくらい見られるものか」だったらどうだろう.
    見られなかった作品ほど配信を終えやすいのなら, 一覧そのものが, これまでに配信された作品全部より, 見られた作品のほうへ寄っていることになる.
    一覧から80本を選ぶ手順をどれほど丁寧にしても, 選ぶ前に一覧から消えた作品は80本の中に入りようがない.
    同じ一覧でも, 知りたい全体が違えば, 偏っているかどうかも違ってくる.

    途中で消えずに残ったものだけが手元に集まって生じるこの偏りを\textbf{生存者バイアス}と呼ぶ.
    成功した経営者の習慣を集めた本に, 早起きや毎朝の運動といった共通点が並んでいるのも同じ形だ.
    同じ習慣を続けていても事業に失敗した人は本に取り上げられないので, その習慣が失敗した人にも同じくらい多かったかどうかは, どこにも書かれていない.

  \subsection{数字は何を測っているか}

    集め方に偏りがなかったとしても, 手元にある数字そのものが何を測っているのか, という問題が残っている.

    ここまで映画80本について調べてきたのは, 1本あたりにならすとどれくらい見られているか, 原作の有無で平均視聴回数に差があるか, だった.
    どちらも視聴回数についての疑問で, 視聴回数を数えているかぎり, ずれはない.

    ところが, 出てきた数を見て「原作のある作品のほうが面白い」「こちらのほうが質が高い」と言いたくなる.
    そこで一歩はみ出してしまう.

    視聴回数が数えているのは, 再生ボタンが押されて動画が始まった回数だ.
    見始めて5分で止めた人の1回も, 最後まで見て二度見返した人の1回も, 同じ「1回」として入っている.
    面白さを直接数えることはできないので, 数えられるものを数えている.

    通販サイトの星の数でも同じことが起きている.
    製品の品質を知りたいと思って星の数を眺めるが, 星1つがついていても, その理由は製品そのものではなく, 配送が予定より遅れたことへの不満かもしれない.
    逆に, 期待がそれほど高くなかったために星5つがついていることもある.
    星の平均を見るときに一人ひとりの評価が均等にならされてしまうことは確かめたが, ならされる前の一つひとつの星にも, 「製品の良し悪し」以外の事情が混ざり込んでいる.

    平均の計算も, 信頼区間も, 検定も, すべて「視聴回数」や「星の数」という手元の数字については正しく行われている.
    問題があるとしたら, 出てきた数字をいつの間にか「面白さ」や「商品の質」そのものの話として解釈してしまうところにある.

    映画の本数をいくら増やしたところで, 視聴回数が「再生が始まった回数」であることに変わりはなく, 知りたいものと測っているものとのずれは縮まらない.

    このずれは, 手元のデータをいくら眺めていても見えてこない.
    数字の計算を見るのではなく, その数字がどのように数えられ, 何を聞いて記録されたものなのかという, データが作られた経緯を確かめるしかない.

  \subsection{選ばれた結果と珍しさの意味}

    集め方にも, 測っているものにも問題がなかったとしても, もう一つ別の場所から数字の解釈を狂わせる問題が入り込む.
    計算された結果の選び方だ.

    ここに1枚のコインがあるとする.
    このコインを100回投げてみたところ, 表が62回出た.
    表と裏が半分ずつ出るふつうのコインなら, 表が出る回数の期待値は50回だ.
    検定の枠組みに沿って, 「このコインは表と裏が同じ確率で出る」という帰無仮説を立て, 有意水準5\%で検定してみる.
    計算される$p$値はおよそ0.02となり, 0.05を下回る.
    帰無仮説は棄却され, 表と裏の出方に偏りがある, つまり「このコインには細工がありそうだ」という結論になる.
    前章までの話のとおりで, 筋の通った結論だ.

    実は, 手元にはコインが最初から20枚あった.
    その20枚のコインをすべて100回ずつ投げ, 20回分の検定を行った.
    残りの19枚は, ひとつも有意にはならなかった (\cref{fig:ch07-coins}).
    そして, $p$値が0.05を下回ったたった1枚のコインだけを, 先ほどのように取り出して見せたのだ.
    1枚1枚の計算には何の誤りもない.

    \begin{figure}[htbp]
      \centering
      % 第7章 図7-2: 細工のないコイン20枚を100回ずつ投げたときの表の回数 (例示用のデータ)
% 19枚は二項分布 (100回, 0.5) から 40〜60回の範囲で乱数生成 (seed 20260927), 13枚目を62回に固定
% 両側の二項検定で有意水準5%を下回るのは 39回以下か61回以上 (60回で p=0.057, 61回で p=0.035, 62回で p=0.021)
\begin{tikzpicture}[font=\footnotesize]
  \begin{axis}[
    width=0.8\linewidth, height=5.2cm,
    xmin=0.2, xmax=20.8, ymin=30, ymax=70,
    xtick={1,5,10,15,20}, ytick={30,40,50,60,70},
    xlabel={コインの番号}, ylabel={表の回数},
    axis x line*=bottom, axis y line*=left,
    clip=true,
  ]
    % 有意になる範囲
    \fill[black!12] (axis cs:0.2,60.5) rectangle (axis cs:20.8,70);
    \fill[black!12] (axis cs:0.2,30) rectangle (axis cs:20.8,39.5);
    \draw[densely dashed, black!60] (axis cs:0.2,50) -- (axis cs:20.8,50);
    \addplot[only marks, mark=*, mark size=2.3pt, mark options={fill=white, draw=black}] coordinates {
      (1,48) (2,43) (3,43) (4,48) (5,47) (6,52) (7,55) (8,50) (9,56) (10,42)
      (11,43) (12,50) (14,53) (15,55) (16,42) (17,51) (18,48) (19,57) (20,45)};
    \addplot[only marks, mark=*, mark size=2.8pt, mark options={fill=black, draw=black}] coordinates {(13,62)};
    \node[font=\footnotesize, anchor=south] at (axis cs:10.5,63.5) {有意になる範囲 (61回以上)};
    \node[font=\footnotesize, anchor=north] at (axis cs:10.5,36.5) {有意になる範囲 (39回以下)};
  \end{axis}
\end{tikzpicture}

      \caption{コイン20枚をそれぞれ100回投げたときの, 表の出た回数. 例示用に作ったデータで, 13枚目 (黒い点) を62回としている. 横軸はコインの番号, 縦軸は表の回数. 灰色の範囲は, 表と裏が同じ確率で出るという帰無仮説を有意水準5\%で検定したときに有意になる回数 (39回以下と61回以上), 破線はその帰無仮説のもとでの期待値50回を示す.}
      \label{fig:ch07-coins}
    \end{figure}

    しかし, このコインには本当に細工があると言っていいのだろうか.
    有意水準5\%という基準を思い出してみよう.
    これは, 「細工のないふつうのコインであっても, 20回に1回ほどは偶然で有意になってしまう」という誤りを認める約束だった.
    ということは, まったく細工のない公平なコインを20枚集めてきて同じ実験をすれば, そのうちのどれか1枚くらいは, 偶然だけで有意になってしまっても何の不思議もない.

    細工のないコインを20枚投げたとき, 少なくとも1枚が偶然有意になってしまう確率は, 次のように計算できる.
    有意水準どおりに, 1枚が偶然有意になる確率を5\%と置いてみる.
    コイン1枚が有意にならない確率は $0.95$ だから, 20枚すべてが有意にならない確率は $0.95^{20}$ になる.
    その裏返しとして, 20枚のうちどれか1枚でも有意になる確率は,
    \[
      1 - 0.95^{20} \approx 0.64
    \]
    とおよそ64\%になる.
    3回に2回ほどの割合で, 細工など何もないふつうのコインの集まりから「有意な偏り」が勝手に生み出されてしまうわけだ.

    1枚だけを投げて有意になるのは, 確かに珍しい出来事だ.
    しかし, 20枚投げてその中から有意になったものだけを拾い上げるなら, それは少しも珍しいことではない.
    その1枚について計算された$p$値の数字は少しも変わっていないのに, 「珍しいことが起きた」という言葉の重みがまるで違ってしまう.

    このように, 望ましい結果が出るまで条件や対象を変えて分析を繰り返し, 有意になったものだけを取り出す行為を\textbf{$p$ハッキング}と呼ぶ.
    そうならないように, データを見る前に, 「どの比較を行うのか」「何を試すのか」をあらかじめ決めておく.
    結果を見てから都合よく動かさないために有意水準を事前に決めておいたのと同じように, 分析の対象そのものも事前に固定しておく必要がある.
    そして, 試したものは, 有意にならなかったものも含めて全部を示す.
    データを眺めていて傾向に気づくのは悪いことではなく, 確かめてみる値打ちのある見当が手に入ったということだ.
    ただ, 気づいたのと同じデータでその珍しさを測れば, 20枚のコインから有意な1枚を拾うのと変わらない.
    確かめるには, 新しくデータを集めて測り直す.

    選んだことで数字の意味が変わるのは, 検定を何度も試したときだけではない.
    ある学校で, 試験の成績が下のほうだった生徒たちを選んで, 特別に指導したとする.
    次の試験では, その生徒たちの平均点が上がった.
    指導の効果が出たように見える.
    だが, 選んだのは一回の試験で下のほうにいた生徒だ.
    標本を取り直せば平均が動いたように, 一人の点数も, 同じ生徒が試験を受け直せば上下する.
    下のほうにいた生徒の中には, たまたまその回の出来が悪かった生徒が混ざっている.
    そうした生徒は, 指導がなくても, 次の回にはふだんの点数のあたりに戻ってくる.

    そのため, 選んだ生徒たちの平均点は, 指導をしなくても前の回より上がりやすい.
    下のほうから選んで測り直すと, ばらつきだけで戻る分が, 指導の効果のように見えてしまうのだ.
    このように, 一回の結果で選んだ集まりを測り直したとき, その平均が全体の平均のほうへ戻っていくことを\textbf{平均への回帰}と呼ぶ.
    20枚のコインと同じで, 点数の計算はどこも間違っていないのに, 誰を選んで測り直したかで, 上がった点数の意味が変わってしまう.
    指導の効果を確かめたいなら, 薬の治験と同じように, 選んだ生徒をくじで二つに分け, 片方だけを指導すればよい.
    どちらの側にも, たまたまその回が悪かった生徒が同じように混ざるので, 戻る分は両方に同じように出るはずだ.
    それでも二つの平均点に差が残れば, その差は指導によるものと考えてよく, 偶然のばらつきで説明できるかどうかは, 前の章までと同じ手順で確かめられる.

  \subsection{繰り返しているものの正体}

    この本では最初から, 一回きりのデータにとらわれず, 「もし同じ調査を何度もやり直したら」という想像を土台にして進んできた.
    標本を取り直せば平均が散らばることを確かめ, 同じ手順で作った区間には100回に95回ほど母平均が入ると言い, 差が本当にないときに検定を繰り返せば有意になるのは20回に1回ほどだと約束した.

    しかし, やり直すといっても, やり直しているものは一つではなかった.
    区間の約束でやり直しているのは, 全体から標本を選ぶところだ.
    全体から偏りなく選ぶことを何度も繰り返すと考えるから, 「100回に95回」は全体の平均についての話になる.
    くじで分けたときにやり直しているのは, 目の前に集まった人たちを二つに分けるところだ.
    分け直すたびに顔ぶれは入れ替わるが, 人数が十分に多ければ, 薬や指導のほかの点では, どちらの側にも同じような人が入る.
    だから, 残った差を薬や指導の効果だと言える.
    偏りなく集めれば, 結果を全体の話として言える.
    くじで分ければ, 差を原因の話として言える.
    この二つは別のことだ.

    さきほどの生徒たちは, 試験の下のほうから選んだのだから, 学校全体から偏りなく集めた集まりではない.
    それでも, くじで分ければ, 選ばれた生徒たちについては指導の効果があったかどうかを確かめられる.
    ただ, 同じ指導が学校のほかの生徒にも同じように効くかどうかまでは言えない.
    そこまで言うには, 集め方のほうが要る.
    治験も同じだ.
    くじが保証してくれるのは, 治験に参加した人たちの中で薬が効いたかどうかまでだ.
    参加していない人にも同じように効くかどうかは, 参加した人たちがどう集まったかにかかっている.

    100回に95回, あるいは20回に1回という確率が何を意味しているのかは, 自分が何を繰り返すことにしているかによって決まる.
    それは数式に数字を入れて計算することだけでなく, どこからデータを集めてどう分け, その数字が何を測っていて, 何を試してどの結果を見せているのかという, 調査や分析の進め方そのものを含んでいる.

  \subsection{条件を確かめて語る}

    映画80本の分析に戻って, 手元にある結果をもう一度見直してみよう.

    手元のデータから導き出された事実は, 筋道としてははっきりしている.
    80本の平均視聴回数は $209.4$ 万回で, 母平均の95\%信頼区間はおよそ156万回から263万回だった.
    そして原作のある作品とない作品とでは平均して120万回ほどの差があり, その差は偶然のばらつきだけでは説明しにくいものだった.

    しかし, この結果を説明するには, 計算の正しさを示すだけでは足りない.
    その数字がどんな前提の上に立っているのかを, 自分の手で確かめておく必要がある.

    まず確かめるべきなのは, この80本の結果をどの全体の話として言うのか, という点だ.
    80本を選んだ一覧には, 見られずに配信を終えた作品が載っていなかった.
    一覧から偏りなく80本を選んでいれば, 計算した区間は, いま配信されている作品全体の平均についての区間だと言える.
    だが, これまでに配信された作品全部の話にしようとすると, 消えた作品が抜けている分, 80本全体の平均は高めに出ているかもしれないし, 原作の有無による差も違っていたかもしれない.

    あわせて確かめるべきなのは, 原作の有無という比較が, 最初からそれだけを調べるつもりで選ばれたのか, という点だ.
    もし, 映画のジャンルや公開月, 上映時間や主演俳優の年齢など, 何十通りもの切り口で視聴回数を手当たり次第に比べ, その中でたまたま$p$値が0.05を下回ったのが「原作の有無」だったのだとしたらどうだろうか.
    それは20枚のコインを投げて有意になった1枚だけを取り出したのとまったく同じで, 偶然のばらつきを拾い上げたにすぎないかもしれない.
    分析を始める前から「原作の有無で差があるか」という一点を確かめようとしていたと分かってはじめて, $p$値の小ささはそのままの重みを持つ.

    いま配信されている作品の一覧から偏りなく選んだ80本なら, 1本あたりの平均視聴回数は209.4万回で, 配信中の作品全体の平均はおよそ156万回から263万回と見積もってよい.
    原作の有無を最初から比べるつもりだったのなら, 原作のある作品のほうが平均して120万回ほど多く見られているという差も, 偶然のばらつきだけでは説明しにくい.
    配信を終えた作品まで含めた, これまでに配信された作品全部の話や, 原作があれば視聴回数が増えると言い切ることは, この80本からはできない.
    それでも, この一覧から集め, この比較を先に決めて分析したのなら, ここまではっきり言える.

  \begin{mybox}{ここまででわかったこと}
    四段階の四つ目「仮定を確かめて判断する」.

    \textbf{この章で言えるようになったこと}: これまでに配信された作品全部の話をするなら, 見られずに配信を終えた作品は一覧に残らないので, 80本は見られた作品のほうへ寄っているかもしれない.
    20枚から有意な1枚を選んだコインは, $p$値が約 $0.02$ でも珍しくない.
    下のほうから選んだ生徒の点数は, 指導がなくても測り直せば上がりやすい.

    \textbf{まだ言えないこと}: この80本の視聴回数だけでは, 原作をつければ作品がどれだけ面白くなるかまでは言えない.

    \textbf{前の章の見方がどう変わったか}: 比べる二つの集まりの条件だけでなく, どこから集め, 何を数え, どの結果を選んだかも確かめる必要があった.
    「同じ手順を繰り返す」には, 計算の前後でしていたことまで入っていた.
  \end{mybox}
